% Part: incompleteness
% Chapter: arithmetization-syntax
% Section: proofs-in-ax

\documentclass[../../../include/open-logic-section]{subfiles}

\begin{document}

\olfileid{inc}{art}{pax}
\olsection{સ્વયંસિદ્ધિમૂલક \usetoken{P}{derivation}}

\begin{explain}
સ્વયંસિદ્ધિમૂલક !!{derivation}sનું
અંકગણિતીકરણ કરવા માટે આપણે
!!{derivation}sને સંખ્યાઓ વડે નિરૂપવાં
પડશે. !!{derivation}s તો માત્ર
!!{formula}sની શ્રેણીઓ છે. તેથી
સ્પષ્ટ રીત એ છે કે દરેક
!!{derivation}ને તેમાંનાં !!{formula}sના
સંકેતોની શ્રેણીના સંકેત વડે નિરૂપવું.
\end{explain}

\begin{defn}
$\delta$ એ !!{formula}s
$!A_1$, \dots,~$!A_n$થી બનેલી
સ્વયંસિદ્ધિમૂલક !!{derivation} હોય,
તો $\Gn{\delta}$ આ છે:
\[
\tuple{\Gn{!A_1}, \dots, \Gn{!A_n}}.
\]
\end{defn}

\begin{ex}
  આ અત્યંત સરળ !!{derivation} જુઓ:
  \begin{derivation}
  1. & $!B \lif (!B \lor !A)$ \\
  2. & $(!B \lif (!B \lor !A)) \lif (!A  \lif (!B \lif (!B \lor !A)))$\\
  3. & $!A  \lif (!B \lif (!B \lor !A))$
  \end{derivation}
  આ !!{derivation}ની ગડેલ-સંખ્યા
  નીચે પ્રમાણે થાય:
  \begin{align*}
  \openTuple\, 
    & \Gn{!B \lif (!B \lor !A)}, \\
    &\Gn{(!B \lif (!B \lor !A)) \lif (!A  \lif (!B \lif (!B \lor !A)))},\\
    & \Gn{!A  \lif (!B \lif (!B \lor !A))} \,\closeTuple.
  \end{align*}
\end{ex}

\begin{explain}
!!{derivation}sનું નિરૂપણ નક્કી કર્યા પછી
આવા !!{derivation}s પર આદિમ
પુનરાવર્તી રીતે ક્રિયા કરી શકાય
અને તેમના આવશ્યક ગુણધર્મો તથા
સંબંધો એ જ રીતે વ્યક્ત કરી શકાય
તે પણ બતાવવું પડે. કેટલીક ક્રિયાઓ
સરળ છે: દાખલા તરીકે,
!!a{derivation}ની ગડેલ-સંખ્યા~$d$
હોય, તો $(d)_{\len{d}-1}$ તેના અંતિમ
!!{formula}ની ગડેલ-સંખ્યા આપે છે.
બીજી કેટલીક ક્રિયાઓ ઘણી કઠિન છે.
તેમને કેવી રીતે સંભાળવી તેનો ઓછામાં
ઓછો રૂપરેખાત્મક ખ્યાલ આપીશું.
હેતુ એ બતાવવાનો છે કે
``$\delta$ એ~$\Gamma$માંથી~$!A$ની
!!a{derivation} છે'' એ સંબંધ
$\delta$ અને~$!A$ની ગડેલ-સંખ્યાઓનો
આદિમ પુનરાવર્તી સંબંધ છે.
\end{explain}

\begin{prop}
\ollabel{prop:followsby}
નીચેના સંબંધો આદિમ પુનરાવર્તી છે:
\begin{enumerate}
\item $!A$ સ્વયંસિદ્ધિ છે.
\item $\delta$ની $i$મી પંક્તિ
  મોડસ પોનેન્સ વડે પ્રમાણિત છે.
\item $\delta$ની $i$મી પંક્તિ
  \QR{} વડે પ્રમાણિત છે.
\item $\delta$ યોગ્ય !!{derivation} છે.
\end{enumerate}
\end{prop}

\begin{proof}
!!{formula}sની ગડેલ-સંખ્યાઓ અને
!!{derivation}sની ગડેલ-સંખ્યાઓ વચ્ચેના
અનુરૂપ સંબંધો આદિમ પુનરાવર્તી છે
તે બતાવવું છે.
\begin{enumerate}
\item સ્વયંસિદ્ધિ-યોજનાઓની સીમિત યાદી
  આપેલી છે. $!A$ કોઈ યોજના દર્શાવેલા
  સ્વરૂપનું હોય તો સ્વયંસિદ્ધિ છે.
  યોજનાઓની યાદી સીમિત હોવાથી,
  દરેક યોજના માટે $!A$ એ સ્વરૂપનું
  છે કે નહીં તે આદિમ પુનરાવર્તી રીતે
  ચકાસી શકાય તે બતાવવું પૂરતું છે.
  દાખલા તરીકે આ યોજના લો:
  \[
  !B \lif (!C \lif !B).
  \]
  એવાં !!{formula}s $!B$ અને $!C$
  હોય કે `$($' સાથે $!B$, પછી
  `$\lif$', પછી `$($', પછી $!C$,
  પછી `$\lif$', પછી $!B$ અને
  છેલ્લે~`$))$' જોડવાથી $!A$ મળે,
  તો $!A$ આ સ્વયંસિદ્ધિ-યોજનાનું
  દૃષ્ટાંત છે. $!A$ની ગડેલ-સંખ્યા
  $n$ માટે અનુરૂપ ગુણધર્મ ચકાસી
  શકીએ છીએ, કારણ કે શ્રેણીઓનું
  જોડાણ આદિમ પુનરાવર્તી છે. વધુમાં,
  સંબંધ સત્ય હોય ત્યારે $!B$ અને
  $!C$ બંને~$!A$નાં ઉપ-!!{formula}s
  છે; તેથી $!B$ અને $!C$ની ગડેલ-સંખ્યાઓ
  $!A$ની ગડેલ-સંખ્યા કરતાં નાની
  છે. આમ વ્યાખ્યાયિત કરીએ:
  \begin{multline*}
  \fn{IsAx}_{!B \lif (!C \lif !B)}(n) \defiff \bexists{b< n}{
    \bexists{c < n}{(\fn{Sent}(b) \land \fn{Sent}(c) \land {}}}\\ n =
  \Gn{(} \concat b \concat \Gn{\lif} \concat \Gn{(} \concat c \concat
  \Gn{\lif} \concat b \concat \Gn{))}).
  \end{multline*}
  દરેક સ્વયંસિદ્ધિ-યોજના માટે આવી
  વ્યાખ્યા હોય, તો તેમનો વિકલ્પયોગ
  $\fn{IsAx}(n)$ ગુણધર્મ વ્યાખ્યાયિત
  કરે છે: ``$n$ સ્વયંસિદ્ધિની
  ગડેલ-સંખ્યા છે.''
\item $\delta$ની $i$મી પંક્તિ
  મોડસ પોનેન્સથી ત્યારે અને માત્ર
  ત્યારે પ્રમાણિત છે જ્યારે એવી
  પંક્તિઓ $j$ અને $k < i$ હોય કે
  પંક્તિ~$j$ પરનું !!{sentence}
  કોઈ સૂત્ર~$!A$ હોય, પંક્તિ~$k$
  પરનું વાક્ય $!A \lif !B$ હોય
  અને પંક્તિ~$i$ પરનું વાક્ય~$!B$ હોય.
  \begin{multline*}
    \fn{MP}(d, i) \defiff \bexists{j < i}{\bexists{k < i}{}}\\
    (d)_k = \Gn{(} \concat (d)_j
      \concat \Gn{\lif} \concat (d)_i \concat \Gn{)}
  \end{multline*}
  સીમિત પરિમાણન, શ્રેણીઓનું જોડાણ
  અને $=$ આદિમ પુનરાવર્તી હોવાથી
  આ સંબંધ પણ આદિમ પુનરાવર્તી છે.
\item $\delta$ની કોઈ પંક્તિ
  \QR{} વડે પ્રમાણિત છે જો તે
  $!B \lif \lforall[x][!A(x)]$
  સ્વરૂપની હોય, કોઈ અગાઉની પંક્તિ
  કોઈ !!{constant}~$c$ માટે
  $!B \lif !A(c)$ હોય અને $c$
  એ~$!B$માં ન આવતો હોય. આ ત્યારે
  અને માત્ર ત્યારે બને જ્યારે
  \begin{enumerate}
  \item કોઈ !!a{sentence}~$!B$ હોય,
  \item મુક્ત રીતે માત્ર ચલ~$x$
    ધરાવતું !!a{formula}~$!A(x)$ હોય,
  \item પંક્તિ~$i$ પર
    $!B \lif \lforall[x][!A(x)]$ હોય,
  \item કોઈ સ્થિર પદ~$c$ માટે અગાઉની
    પંક્તિ $j < i$ પર
    $!B \lif \Subst{!A}{c}{x}$ હોય,
  \item અને એ~$!B$માં ન આવતું હોય.
  \end{enumerate}
  આ બધી શરતો આદિમ પુનરાવર્તી રીતે
  ચકાસી શકાય છે: $!B$, $!A(x)$
  અને $x$ની ગડેલ-સંખ્યાઓ પંક્તિ~$i$
  પરના સૂત્રની ગડેલ-સંખ્યા કરતાં
  નાની છે, અને $c$ની સંખ્યા
  પંક્તિ~$j$ પરના સૂત્રની ગડેલ-સંખ્યા
  કરતાં નાની છે:
  \begin{multline*}
    \fn{QR}_1(d, i) \defiff \bexists{j<i}{\bexists{b < (d)_i}{\bexists{x <
        (d)_i}{\bexists{a < (d)_i}{\bexists{c < (d)_j}{(}}}}}
    \\ \fn{Var}(x) \land \fn{Const}(c) \land {} \\
    (d)_i = \Gn{(} \concat
    b \concat \Gn{\lif} \concat \Gn{\lforall} \concat x \concat a
    \concat \Gn{)} \land {}\\
    (d)_j = \Gn{(} \concat b \concat
    \Gn{\lif} \concat \fn{Subst}(a,c,x) \concat \Gn{)} \land {}\\ \fn{Sent}(b)
    \land \fn{Sent}(\fn{Subst}(a,c,x)) \land {} \bforall{k <
      \len{b}}{(b)_k \neq (c)_0})
  \end{multline*}
  \sourcecorrection{OLINC-018}{સ્થિર અંગ્રેજી
  મૂળના આ સૂત્રમાં અગાઉની પંક્તિનો
  સૂચક~$j$ તેની મર્યાદા
  $(d)_j$માં અને પછીના સમીકરણમાં
  વપરાય છે, પરંતુ તેને પરિમાણિત
  કરવામાં આવ્યો નથી. અહીં
  $j<i$ની સીમિત અસ્તિત્વ-શરત
  બહારથી ઉમેરી છે. મૂળ ગદ્યમાં
  સ્થિર પદ~$c$ને એક જગ્યાએ
  ભૂલથી~$a$ લખ્યું છે;
  અહીં~$c$ જ રાખ્યું છે.}
  અહીં $c$ અને $x$ અનુક્રમે
  સ્થિર પદ અને ચલની ગડેલ-સંખ્યાઓ છે;
  માત્ર તેમના ચિહ્ન-સંકેતો નથી.
  $!A(x)$નું એકમાત્ર મુક્ત ચલ~$x$
  છે તે ચકાસવા માટે
  $\Subst{!A(x)}{c}{x}$
  !!a{sentence} છે કે નહીં તે ચકાસીએ.
  $!B$નું દરેક ચિહ્ન~$c$થી જુદું છે
  એવી શરતથી $c$ એ~$!B$માં નથી
  તેની ખાતરી કરીએ.

  \QR{}નું બીજું સ્વરૂપ અભ્યાસપ્રશ્ન
  તરીકે મૂકીએ છીએ.
\item $d$ યોગ્ય !!{derivation}ની
  ગડેલ-સંખ્યા ત્યારે અને માત્ર ત્યારે
  છે જ્યારે તેની દરેક પંક્તિ કાં તો
  સ્વયંસિદ્ધિ હોય અથવા મોડસ પોનેન્સ
  કે~\QR{} વડે પ્રમાણિત હોય. તેથી:
  \[
  \fn{Deriv}(d) \defiff \bforall{i < \len{d}}{(\fn{IsAx}((d)_i) \lor
    \fn{MP}(d,i) \lor \fn{QR}(d, i))}
  \]
\end{enumerate}
\end{proof}

\begin{prob}
\olref[inc][art][pax]{prop:followsby}ની
જેમ નીચેના સંબંધો વ્યાખ્યાયિત કરો:
\begin{enumerate}
\item $\fn{IsAx}_{!A \lif (!B \lif (!A \land !B))}(n)$,
\item $\fn{IsAx}_{\lforall[x][!A(x)] \lif !A(t)}(n)$,
\item $\fn{QR}_{2}(d, i)$
  (\QR{}ના બીજા સ્વરૂપ માટે).
\end{enumerate}
\end{prob}

\begin{prop}
$\Gamma$ એ !!{sentence}sનો આદિમ
પુનરાવર્તી ગણ હોય એમ ધારો.
ત્યારે $\Prf[\Gamma](x, y)$ આદિમ
પુનરાવર્તી છે: અહીં~$x$ એ
મૂળ નિષ્કર્ષ~$!A$ માટે
$\Gamma$માંથી મળતાં સાન્ત
ધારણાવાક્યોને પૂર્વપક્ષ બનાવીને
રચાયેલા શરતી વાક્યની બંધ
!!a{derivation}~$\delta$નો સંકેત છે
અને $y$ એ મૂળ નિષ્કર્ષ~$!A$ની
ગડેલ-સંખ્યા છે.
\sourcecorrection{OLINC-019}{સ્થિર અંગ્રેજી
મૂળ આ વિધાનમાં~$x$ને
$\Gamma$માંથી~$!A$ના સીધા
નિગમનનો સંકેત કહે છે,
પરંતુ તેની નીચેની વ્યાખ્યામાં~$x$
ખાલી ધારણાઓમાંથી સાન્ત
$\Gamma$-વાક્યોના શરતીકરણના
નિગમનનો સંકેત છે.
આ બે સંકેત-સંબંધો જોડીદીઠ સમાન
નથી. અહીં વિધાનને નીચે ખરેખર
વ્યાખ્યાયિત થયેલા સંબંધ સાથે
મેળમાં કર્યું છે; મૂળ નિષ્કર્ષની
સાધ્યતા સાથેનો સંબંધ નિગમન
પ્રમેય આપે છે.}
\end{prop}

\begin{proof}
``$y \in \Gamma$'' આદિમ પુનરાવર્તી
વિધાન~$R_\Gamma(y)$ વડે આપવામાં
આવ્યું હોય એમ ધારો. આપણે બતાવવું છે
કે $\Prf[\Gamma](x, y)$ સંબંધ આદિમ
પુનરાવર્તી છે. અંગ્રેજી મૂળની
શરૂઆતમાં $\Prf[\Gamma](x, y)$નો
આશય એવો આપ્યો છે કે
$y$ એ !!a{sentence}~$!A$ની
ગડેલ-સંખ્યા છે અને $x$ એ
$\Gamma$માંથી~$!A$ની
!!a{derivation}નો સંકેત છે. નીચેની
રચના જોકે પુરાવા-સંકેત માટે બંધ નિગમનનો
સંકેત લે છે; આ ભેદ ઉપર નોંધ્યો છે.

અગાઉના વિધાન પ્રમાણે,
$\fn{Deriv}(x)$ આદિમ પુનરાવર્તી છે:
તે ત્યારે અને માત્ર ત્યારે સત્ય છે
જ્યારે $x$ યોગ્ય
!!{derivation}~$\delta$નો સંકેત હોય.
જોકે તે વ્યાખ્યામાં !!{derivation}ની પંક્તિઓને
પ્રમાણિત કરવાની વધારાની રીત તરીકે
ધારણાગણ લેવાયો નથી. \QR{} વડે
પંક્તિ પ્રમાણિત છે કે નહીં તેની
આપણી આદિમ પુનરાવર્તી ચકાસણીમાં
સ્થિર પદ~$c$ એ~$\Gamma$માં આવવું
જોઈએ નહીં એવી શરત પણ સામેલ નથી.
વ્યાખ્યામાં સીધો~$\Gamma$ લઈએ તેવો
ફેરફાર શક્ય છે; પરંતુ
$\fn{Deriv}$ અને નિગમન પ્રમેય
વાપરવા સહેલા છે.
$\Gamma \Proves !A$ ત્યારે અને માત્ર
ત્યારે જ્યારે~$\Gamma$માંનાં
!!{sentence}sની કોઈ સાન્ત યાદી
$!B_1$, \dots, $!B_n \in
\Gamma$ એવી હોય કે
$\{!B_1, \dots, !B_n\} \Proves !A$.
નિગમન પ્રમેય પ્રમાણે આ ત્યારે
સત્ય છે જ્યારે
$\Proves (!B_1 \lif (!B_2 \lif
\cdots (!B_n \lif !A)\cdots))$.
ગડેલ-સંખ્યા~$z$ ધરાવતું
!!a{sentence} આ સ્વરૂપનું છે કે
નહીં તે આદિમ પુનરાવર્તી રીતે
ચકાસી શકાય છે. તેથી $x$ને
ગડેલ-સંખ્યા~$y$ ધરાવતા
!!{sentence}ની \emph{$\Gamma$માંથી}
!!a{derivation}નો સંકેત લેવા બદલે,
આ~$x$ને ઉપરના સ્વરૂપના નેસ્ટેડ
શરતી વાક્યની~$\emptyset$માંથી
!!a{derivation}નો સંકેત લઈએ છીએ.

પહેલાં, !!{sentence}sની શ્રેણી હોય,
તો એ બધાં વાક્યો પૂર્વપક્ષમાં
અને આપેલું !!{sentence} ઉત્તરપક્ષમાં
રાખીને આદિમ પુનરાવર્તી રીતે
શરતી વાક્ય બનાવી શકીએ:
\begin{align*}
  \fn{hCond}(s, y, 0) & = y \\
  \fn{hCond}(s, y, n+1) & = \Gn{(} \concat (s)_{n} \concat \Gn{\lif}
  \concat \fn{hCond}(s, y, n) \concat \Gn{)}\\
  \fn{Cond}(s, y) & = \fn{hCond}(s, y, \len{s})\\
  \intertext{તેથી $\Prf[\Gamma](x, y)$ને આ પ્રમાણે વ્યાખ્યાયિત કરી શકીએ:}
  \Prf[\Gamma](x, y) & \defiff \bexists{s < \fn{sequenceBound}(x,x)}{(} \\
  &\qquad (x)_{\len{x}-1} = \fn{Cond}(s,y) \land {} \\
  &\qquad\bforall{i<\len{s}}{R_\Gamma((s)_i)} \land {} \\
  &\qquad\fn{Deriv}(x)).
\end{align*}
\sourcecorrection{OLINC-020}{સ્થિર અંગ્રેજી
મૂળમાં $\fn{hCond}$ના પુનરાવર્તનના
પગલે ત્રણ દલીલોવાળું
$\fn{Cond}(s,y,n)$ લખ્યું છે,
જ્યારે $\fn{Cond}$ની વ્યાખ્યા
નીચે માત્ર બે દલીલો માટે છે.
અહીં પુનરાવર્તનમાં યોગ્ય
$\fn{hCond}(s,y,n)$ રાખ્યું છે.}
\sourcecorrection{OLINC-021}{સ્થિર અંગ્રેજી
મૂળના અંતિમ સૂત્રમાં
$(s)_i \in \Gamma$ અનૌપચારિક
સભ્યત્વ-લેખન છે; ઉપરથી આપેલા
આદિમ પુનરાવર્તી વિધાન
$R_\Gamma((s)_i)$ વડે જ તેની
સીમિત ચકાસણી વ્યક્ત થાય છે.
અહીં એ જ વિધાન વાપર્યું છે.}
$s$ માટેની મર્યાદા એ હકીકત પરથી
આવે છે કે દરેક $(s)_i$ એ
!!{derivation}ની છેલ્લી પંક્તિના
કોઈ ઉપ-!!{formula}ની ગડેલ-સંખ્યા
છે; એટલે તે $(x)_{\len{x}-1}$થી
નાની છે. $!B \in \Gamma$
પૂર્વપક્ષોની સંખ્યા, એટલે~$s$ની
લંબાઈ, $x$ની છેલ્લી પંક્તિની
લંબાઈ કરતાં નાની છે.
\end{proof}

\end{document}
